\section{Related work}
\label{sec:related}

\paragraph{Agent improvement loops.} Karpathy's \emph{autoresearch}~\cite{autoresearch} popularized giving a
coding agent a fixed training script, a metric and a loop of edit, run, keep or discard. Community forks extend
it to other targets~\cite{uditgoenka,tbelbek}. These are recipes; none reports a controlled comparison against an
agent without the loop, or measures what the loop ships beyond its own metric. hill-climb is one such recipe,
packaged as an agent skill, and this report is that comparison.

\paragraph{Optimization benchmarks for agents.} GSO~\cite{gso} and PerfAgent~\cite{perfagent} score repository-level
performance patches against hidden correctness tests and human speedups, and report reward hacking such as caching
inside timing loops. They rank models. We hold the model fixed and vary only the process instructions, and we add
two measures they do not report: whether the agent's own claim matches an outside measurement, and whether the
final version keeps behavior that only a hidden gate checks.

\paragraph{LLM judges.} Pairwise LLM judging is known to carry position and self-preference biases~\cite{mtbench}.
For copy, LOLA~\cite{lola} found that prompt-based LLMs pick the better of two Upworthy headlines no more than
65\% of the time. We test the specific protocol hill-climb uses for taste work (blind, per-judge coin flips,
three fresh judges, strict majority) on the same archive~\cite{upworthy}, and ask a question those studies do not:
does a panel's agreement tell you it is right?
